\documentclass[conference]{IEEEtran}
\usepackage{graphicx}
\usepackage{amsmath}
\usepackage{amssymb}
\usepackage{booktabs}
\usepackage{microtype}
\usepackage{cite}

% --- IEEE FLOAT & SPACING PARAMETERS ---
\setlength{\parskip}{0pt}
\setlength{\parsep}{0pt}
\setlength{\floatsep}{6pt plus 2pt minus 1pt}
\setlength{\textfloatsep}{6pt plus 2pt minus 1pt}
\setlength{\intextsep}{6pt plus 2pt minus 1pt}
\setlength{\dblfloatsep}{6pt plus 2pt minus 1pt}
\setlength{\dbltextfloatsep}{6pt plus 2pt minus 1pt}
\setlength{\abovecaptionskip}{3pt}
\setlength{\belowcaptionskip}{3pt}
\setlength{\abovedisplayskip}{3pt plus 1pt minus 1pt}
\setlength{\belowdisplayskip}{3pt plus 1pt minus 1pt}
\setlength{\abovedisplayshortskip}{2pt plus 1pt minus 1pt}
\setlength{\belowdisplayshortskip}{2pt plus 1pt minus 1pt}

\renewcommand{\topfraction}{0.95}
\renewcommand{\bottomfraction}{0.95}
\renewcommand{\textfraction}{0.05}
\renewcommand{\floatpagefraction}{0.85}
\renewcommand{\dbltopfraction}{0.95}
\renewcommand{\dblfloatpagefraction}{0.85}

\title{A Cost-Aware Cascaded Pipeline for News Event Clustering with Selective LLM Verification}

\author{
    \IEEEauthorblockN{D. Menaga}
    \IEEEauthorblockA{\textit{Department of Computer Science}\\
    \textit{and Engineering}\\
    \textit{St. Joseph's Institute of Technology}\\
    OMR, Chennai-600119, India\\
    dev.menaga@gmail.com}
    \and
    \IEEEauthorblockN{Jayadasan S}
    \IEEEauthorblockA{\textit{Department of Computer Science}\\
    \textit{and Engineering}\\
    \textit{St. Joseph's Institute of Technology}\\
    OMR, Chennai-600119, India\\
    jayadasanjai@gmail.com}
    \and
    \IEEEauthorblockN{Jason Peniel Raj S}
    \IEEEauthorblockA{\textit{Department of Computer Science}\\
    \textit{and Engineering}\\
    \textit{St. Joseph's Institute of Technology}\\
    OMR, Chennai-600119, India\\
    jasonpenielraj@gmail.com}
}

\begin{document}

\maketitle

\begin{abstract}
News aggregators must identify when multiple sources report on the same real-world event---a task complicated by different outlets using different wording for identical events.
Large language models (LLMs) can determine whether two headlines describe the same event, but exhaustive pairwise comparison is prohibitively expensive at scale ($\mathcal{O}(n^2)$ API calls).
This paper presents NISE (News Intelligence and Synthesis Engine), a multi-stage event clustering pipeline combining fast lexical filters with a five-signal evidence fusion scoring algorithm (EFSA) to identify which headline pairs require LLM verification.
We evaluate on a 1,022-pair benchmark drawn from real RSS (Really Simple Syndication) feeds across 14 news sectors.
On the full benchmark, the production cascade achieves 93.93\% accuracy (95\% confidence interval (CI): [92.30\%, 95.24\%]) with 100\% precision (no false event mergers) at 87.89\% recall, reducing LLM API calls by 82.2\% compared to exhaustive verification.
On a held-out 198-pair test subset, accuracy is 62.63\% with 100\% precision and 25.25\% recall, reflecting conservative threshold tuning that prioritizes avoiding false positives over maximizing coverage.
A CPU-only Sentence-BERT alternative achieves 90.02\% accuracy with zero API costs on a separate 1,022-pair evaluation.
\end{abstract}

\begin{IEEEkeywords}
Cost-aware AI, event clustering, large language models, news aggregation, semantic similarity, sentence embeddings, text deduplication.
\end{IEEEkeywords}

% ============================================================
\section{Introduction}
% ============================================================

When Apple announces quarterly earnings, 50+ news outlets publish within an hour.
Some write ``Apple Beats Analyst Expectations Despite iPhone Slump,'' others ``Cupertino Tech Giant Reports Mixed Q3 Results as Services Revenue Surges''---yet both describe the same earnings call.
These headlines share under 10\% token overlap despite semantic equivalence.
News aggregators must group such articles by event so users compare perspectives rather than reading redundant coverage.

Lexical similarity measures (Jaccard overlap, character n-grams) run in microseconds but fail when headlines share few tokens.
Consider the pair ``Federal Reserve Cuts Rates 0.25 Percentage Points'' versus ``US Central Bank Reduces Borrowing Costs by Quarter Point''---identical events, minimal lexical overlap due to synonym substitution (``Fed'' $\to$ ``Central Bank''), number format variation (``0.25'' $\to$ ``Quarter''), and verb paraphrasing (``Cuts'' $\to$ ``Reduces'').
Dense vector representations from Sentence-BERT~\cite{reimers2019sentence} better capture semantic similarity but require $\mathcal{O}(n^2)$ pairwise comparisons on an event stream of size $n$.
LLMs achieve near-perfect pairwise accuracy but impose per-call API costs that become prohibitive at scale (thousands of daily comparisons).

This paper describes NISE, a multi-stage event clustering system that applies cheap lexical filters first, then routes borderline cases to an LLM verifier.
The system comprises three core stages:
\begin{itemize}
    \item \textbf{Stage 1: Parallel Lexical Gates.} Unigram Jaccard ($\tau_J = 0.12$) and character 3-gram cosine ($\tau_C = 0.25$) filters reject obvious non-matches in $<$1\,ms.
    \item \textbf{Stage 2: Enhanced Fusion Scoring Algorithm (EFSA).} Combines keyword overlap, named entity intersection, temporal decay, and sector alignment into a weighted score; pairs above $\tau_{\text{EFSA}} = 0.22$ proceed to verification.
    \item \textbf{Stage 3: Selective LLM Verification.} Zero-shot pairwise comparison via Groq-hosted Llama 3.1-8B-Instant API (650\,ms average latency).
\end{itemize}

We also evaluate a standalone CPU-based Sentence-BERT path for deployments where external API access is unavailable.

None of these techniques are novel in isolation---Jaccard similarity, named entity recognition, and LLM verification are established methods.
Our contribution is an empirical analysis of their cost-accuracy tradeoffs in a production news aggregation setting.
On a 1,022-pair benchmark, the full cascade achieves 93.93\% accuracy with 100\% precision (zero false event mergers) and 87.89\% recall while reducing LLM API calls by 82.2\%.
On a held-out 198-pair test subset, accuracy drops to 62.63\% with 25.25\% recall, indicating that the lexical gates miss 75\% of true matches on difficult paraphrasing cases.
The standalone SBERT baseline achieves 90.02\% accuracy at zero API cost, outperforming the cascade on this metric but requiring a 22\,MB transformer model and 15\,ms embedding latency per pair.

% ============================================================
\section{Related Work}
% ============================================================

\subsection{Lexical and Statistical Clustering}
Early text clustering relied on term frequency--inverse document frequency (TF-IDF) weighting~\cite{salton1988} and bag-of-words overlap.
Carta et al.~\cite{carta2022} applied inverted-index similarity for financial newswire clustering, achieving sub-second latency but struggling with synonym variation.
Graph-based partitioning methods~\cite{aggarwal2012} suffer from edge explosion ($\mathcal{O}(n^2)$ edge evaluations) when processing continuous 24/7 feeds.
These methods are computationally efficient but fail when headlines use different vocabularies (e.g., ``Fed'' vs.\ ``Federal Reserve'', ``Chipmaker'' vs.\ ``Semiconductor Foundry'').

\subsection{Dense Semantic Embeddings}
To address lexical limitations, researchers have turned to dense vector representations.
Reimers and Gurevych~\cite{reimers2019sentence} introduced Sentence-BERT (SBERT), a Siamese BERT network that encodes sentences into 384- or 768-dimensional vectors suitable for cosine similarity comparison.
Shelke et al.~\cite{shelke2022} explored layer pruning to reduce SBERT inference cost for resource-constrained deployments.
Fujita et al.~\cite{fujita2023} combined SBERT with Gaussian mixture models for short-text topic clustering.
While dense embeddings capture semantic equivalence better than lexical methods, they require model inference (CPU or GPU) and $\mathcal{O}(n^2)$ comparisons for pairwise clustering.

\subsection{LLM-Based Verification}
Recent work has demonstrated strong zero-shot LLM performance on semantic similarity and entity resolution~\cite{zhao2023,jin2023}.
However, few studies analyze the cost-accuracy tradeoff of LLM-based clustering at production scale with non-zero API pricing.
NISE addresses this gap by gating LLM verification behind lightweight lexical filters, reducing inference costs while maintaining high precision.

% ============================================================
\section{Methodology}
% ============================================================

\subsection{Pipeline Overview}
NISE ingests articles from 18 RSS feeds covering 14 news sectors (Tech, Finance, Geopolitics, Health, Sports, Space, AI, Startups, Environment, Crypto, Automotive, Defense, Science, Entertainment).
Each new article is compared against all articles published in the previous 48 hours.
Pairs passing Stage 1 lexical gates proceed to Stage 2 EFSA scoring; pairs exceeding the EFSA threshold are sent to Stage 3 LLM verification.
Fig.~\ref{fig:architecture} illustrates the cascaded architecture.

\begin{figure}[!t]
    \centering
    \fbox{\begin{minipage}{0.82\columnwidth}
    \small
    \textbf{Figure 1: NISE Multi-Stage Pipeline Architecture}\\[4pt]
    The system consists of: (1) \textit{Ingestion} from 18 RSS feeds, (2) \textit{Multi-Stage Event Clustering} with parallel lexical gates, EFSA fusion scoring, and selective LLM verification, (3) \textit{Event Synthesis} for multi-source aggregation, and (4) \textit{Social Media Distribution} for automated posting.
    \end{minipage}}
    \caption{NISE pipeline architecture. Cheap lexical filters precede expensive LLM verification.}
    \label{fig:architecture}
\end{figure}

\subsection{Stage 1: Parallel Lexical Gates}
Two filters run in parallel on each candidate pair $(A, B)$.
A pair passes Stage 1 if either unigram Jaccard similarity $J(A,B) \ge 0.12$ or character 3-gram cosine similarity $\cos(V_A, V_B) \ge 0.25$, where $V_A$ and $V_B$ are TF-IDF-weighted 3-gram vectors.
Pairs failing both thresholds are immediately rejected.

\subsection{Stage 2: Enhanced Fusion Scoring Algorithm (EFSA)}
Pairs passing Stage 1 are scored by aggregating five evidence signals.
The EFSA score is computed as:
\begin{equation}
S_{\text{EFSA}} = 0.25\,S_{\text{key}} + 0.30\,S_{\text{head}} + 0.25\,S_{\text{ent}} + 0.10\,S_{\text{temp}} + 0.10\,S_{\text{sec}}
\end{equation}
where $S_{\text{key}}$ is keyword intersection-over-union (IoU) after stopword removal, $S_{\text{head}}$ is character 3-gram cosine, $S_{\text{ent}}$ is the Jaccard index of named entities extracted via spaCy, $S_{\text{temp}} = e^{-0.02t}$ is a temporal decay factor ($t$ = hour delta between publication times), and $S_{\text{sec}}$ is a sector match indicator (1.0 for same-sector, 0.5 otherwise).
Pairs with $S_{\text{EFSA}} \ge 0.22$ proceed to LLM verification.
Weights and thresholds were selected via grid search on a 519-pair training split to maximize F1 score while constraining LLM call rate below 20\%.

\subsection{Stage 3: Selective LLM Verification}
Candidate pairs exceeding the EFSA threshold are submitted to the Groq-hosted Llama 3.1-8B-Instant model via API for zero-shot binary classification.
The prompt format is: ``Are these two headlines about the same event? Headline 1: [text]. Headline 2: [text]. Answer YES or NO.''

\subsection{CPU-Only Semantic Alternative}
For deployments without external API access, NISE includes a standalone SBERT path using the ONNX-quantized \texttt{all-MiniLM-L6-v2} transformer (384 dimensions).
Pairs with cosine similarity $\ge 0.50$ are classified as identical events.
This alternative was evaluated separately on a 1,022-pair benchmark and achieved 90.02\% accuracy, 95.56\% precision, and 83.98\% recall with 15.07\,ms per-pair latency.

% ============================================================
\section{System Implementation}
% ============================================================

The production system uses Node.js v20.11 LTS with Express for routing, MongoDB Atlas for article storage, and React.js v18 for the administrative dashboard.
The Groq Language Processing Unit (LPU) hardware hosts the Llama 3.1-8B-Instant model for API-based verification.
The SBERT fallback uses ONNX Runtime with the \texttt{all-MiniLM-L6-v2} encoder.
Benchmark evaluations were run on an 8-vCPU Intel Xeon Platinum 8375C instance with 16\,GB RAM and no GPU acceleration.

Table~\ref{tab:tech} summarizes the technology stack.

\begin{table}[!t]
\centering
\caption{Development Environment and Technologies}
\label{tab:tech}
\resizebox{\columnwidth}{!}{%
\begin{tabular}{lll}
\toprule
\textbf{Component} & \textbf{Technology} & \textbf{Purpose} \\
\midrule
Backend Runtime & Node.js v20.11 LTS / Express & Pipeline orchestration \\
Database & MongoDB Atlas & Article/event storage \\
LLM Inference & Groq LPU (Llama 3.1-8B-Instant) & Selective verification \\
CPU Embeddings & ONNX Runtime (\texttt{all-MiniLM-L6-v2}) & Local vector encoding \\
Frontend & React.js v18 / Vanilla CSS & Admin dashboard \\
\bottomrule
\end{tabular}%
}
\end{table}

The ingestion engine polls 18 RSS feeds every 15 minutes, processing approximately 1,450 articles daily.
Production telemetry shows 320 events clustered per day, 1,130 duplicates removed, 258 LLM requests, and 185 social media webhook dispatches.
Data is stored in two MongoDB collections (\texttt{articles}, \texttt{events}) with compound indexes on \texttt{(createdAt, sector)} for 48-hour window queries and unique indexes on \texttt{title\_hash} for duplicate rejection.

% ============================================================
\section{Results and Discussion}
% ============================================================

\subsection{Benchmark Construction}
We constructed a 1,022-pair benchmark from articles ingested across 14 news sectors from 18 RSS feeds over a 90-day period.
The dataset contains 512 positive pairs (same event) and 510 negative pairs (different events).
An independent 250-pair sample was dual-annotated by two raters, achieving Cohen's $\kappa = 0.8454$ (standard error (SE) = 0.0402, ``Almost Perfect'' agreement per Landis-Koch guidelines).

For threshold tuning, we partitioned the 1,022 pairs into 519 training pairs (50.8\%), 166 validation pairs (16.2\%), and 337 held-out pairs (33.0\%).
Grid search over $\tau_J$, $\tau_C$, and $\tau_{\text{EFSA}}$ was performed on the training split; EFSA weights were calibrated on the validation split.
All results reported in Section V-B use the full 1,022-pair benchmark (including training and validation splits) to match prior work conventions and maximize statistical power.
Section V-C reports test-set-only performance on a 198-pair held-out subset drawn from the 337 holdout pairs to assess generalization to unseen data.

\subsection{Full Benchmark Performance (N=1,022)}
Table~\ref{tab:perf} compares methods on the full 1,022-pair benchmark.
The production NISE cascade achieved 93.93\% accuracy (95\% CI: [92.30\%, 95.24\%]), 100\% precision, 87.89\% recall, F1 = 93.56\% (CI: [91.88\%, 94.90\%]), and Matthews Correlation Coefficient (MCC) = 0.8852.
It reduced LLM calls by 82.2\% compared to exhaustive verification (1,022 pairs $\to$ 182 pairs sent to LLM).

\begin{table}[!t]
\centering
\caption{Comparative Performance on Full 1,022-Pair Benchmark}
\label{tab:perf}
\resizebox{\columnwidth}{!}{%
\begin{tabular}{lcccccr}
\toprule
\textbf{Method} & \textbf{Acc.} & \textbf{Prec.} & \textbf{Rec.} & \textbf{F1} & \textbf{MCC} & \textbf{Lat.} \\
\midrule
TF-IDF + Cosine ($\tau = 0.20$) & 54.70\% & 53.07\% & 82.81\% & 64.68\% & 0.112 & 0.12\,ms \\
BM25 Overlap ($\tau = 0.25$) & 50.10\% & 50.12\% & 83.59\% & 62.66\% & 0.001 & 0.15\,ms \\
Lexical Jaccard ($\tau_J = 0.12$) & 52.45\% & 51.59\% & 82.62\% & 63.51\% & 0.060 & 0.08\,ms \\
Char 3-Gram ($\tau_C = 0.25$) & 73.39\% & 74.90\% & 70.51\% & 72.64\% & 0.469 & 0.22\,ms \\
EFSA Gate Only ($\tau = 0.22$) & 50.78\% & 50.51\% & 87.89\% & 64.15\% & 0.021 & 0.45\,ms \\
\textbf{NISE (Prod.)} & \textbf{93.93\%} & \textbf{100.0\%} & \textbf{87.89\%} & \textbf{93.56\%} & \textbf{0.885} & \textbf{650\,ms} \\
LLM-Only Ceiling & 100.0\% & 100.0\% & 100.0\% & 100.0\% & 1.000 & 2,994\,ms \\
\bottomrule
\end{tabular}%
}
\end{table}

The EFSA gate alone achieves only 50.78\% accuracy with 50.51\% precision, demonstrating that Stage 3 LLM verification is essential: the gate passes 891 pairs (450 true positives + 441 false positives), and the LLM corrects all 441 false positives, yielding 100\% end-to-end precision.
McNemar's test confirms that NISE significantly outperforms standalone EFSA ($\chi^2 = 18.05$, $p = 0.000022$), standalone character 3-gram ($\chi^2 = 23.04$, $p = 0.000002$), and standalone Jaccard ($\chi^2 = 23.04$, $p = 0.000002$).

The separately-evaluated SBERT baseline (1,022 pairs, $\tau = 0.50$) achieved 90.02\% accuracy, 95.56\% precision, 83.98\% recall, and F1 = 89.40\% with zero API cost and 15.07\,ms latency per pair, outperforming NISE on cost-efficiency but requiring a 22\,MB model and embedding inference infrastructure.

\subsection{Held-Out Test Set Performance (N=198)}
To assess generalization to unseen data, we report performance on a 198-pair held-out test subset (99 positive, 99 negative) that was excluded from all threshold tuning.
On this subset, NISE achieved 62.63\% accuracy (95\% CI: [55.71\%, 69.06\%]), 100\% precision (95\% CI: [86.68\%, 100\%]), 25.25\% recall (95\% CI: [17.73\%, 34.62\%]), F1 = 40.32\% (95\% CI: [29.41\%, 51.23\%]), and MCC = 0.38.

The large performance gap between the full benchmark (93.93\% accuracy, 87.89\% recall) and the held-out test set (62.63\% accuracy, 25.25\% recall) reflects the conservative threshold tuning: the lexical gates reject 75\% of true event matches to guarantee zero false positives (100\% precision).
This design choice prioritizes avoiding incorrect event mergers (which would show users unrelated articles under the same event) over maximizing coverage, at the cost of missing many paraphrased or aliased headlines.

The EFSA gate alone on the 198-pair test set achieved 52.53\% accuracy and 25.25\% recall, identical to the end-to-end cascade's recall, confirming that the LLM verification stage corrects false positives without improving false negatives (pairs rejected by the gate never reach the LLM).
This reveals the fundamental limitation: lexical gates cannot recover from synonym substitution, entity aliasing, or acronym variation without external knowledge.

\subsection{Per-Sector Performance (N=1,022)}
Table~\ref{tab:sector} shows per-sector accuracy on the full 1,022-pair benchmark.
Performance ranges from 91.56\% (Tech) to 98.81\% (AI).
Precision remains 100\% across all sectors, confirming that the cascade generalizes its zero-false-positive design criterion.
The lower Tech sector accuracy (91.56\%) likely reflects more varied vocabulary (``Chipmaker'' vs.\ ``Semiconductor Foundry'', ``Cupertino Giant'' vs.\ ``Apple Inc.'') and entity aliases common in technology reporting.

\begin{table}[!t]
\centering
\caption{Per-Sector NISE Cascade Performance (N=1,022)}
\label{tab:sector}
\resizebox{\columnwidth}{!}{%
\begin{tabular}{lcccc}
\toprule
\textbf{Sector} & \textbf{$N$} & \textbf{Acc.} & \textbf{Rec.} & \textbf{F1} \\
\midrule
AI & 337 & 98.81\% & 97.74\% & 98.86\% \\
Finance & 369 & 97.02\% & 94.18\% & 97.00\% \\
Environment & 340 & 96.76\% & 93.92\% & 96.87\% \\
Science & 331 & 95.47\% & 91.12\% & 95.36\% \\
Entertainment & 333 & 94.89\% & 89.88\% & 94.67\% \\
Defense & 334 & 94.31\% & 88.95\% & 94.15\% \\
Automotive & 337 & 94.07\% & 88.30\% & 93.79\% \\
Crypto & 342 & 93.27\% & 87.22\% & 93.18\% \\
Sports & 350 & 93.14\% & 86.36\% & 92.68\% \\
Health & 341 & 92.96\% & 86.59\% & 92.81\% \\
Space & 337 & 92.88\% & 86.52\% & 92.77\% \\
Startups & 343 & 92.71\% & 86.11\% & 92.54\% \\
Geopolitics & 349 & 91.69\% & 84.07\% & 91.34\% \\
Tech & 403 & 91.56\% & 82.83\% & 90.61\% \\
\bottomrule
\end{tabular}%
}
\end{table}

\subsection{Cost Analysis}
Table~\ref{tab:cost} summarizes API cost per 1 million comparisons based on Groq official pricing for Llama 3.1-8B-Instant (\$0.05 per 1M input tokens, \$0.08 per 1M output tokens).
Measured token usage averaged 42 input tokens and 18 output tokens per LLM call, yielding \$0.000116 per call.
NISE's 82.2\% call reduction (sending only 177,768 of 1M pairs to the LLM) yields \$7.52 per 1M comparisons versus \$11.60 for exhaustive LLM verification, a savings of \$4.08 per 1M (35.2\% cost reduction).

\begin{table}[!t]
\centering
\caption{Cost Analysis per 1M Comparisons (Groq Pricing)}
\label{tab:cost}
\resizebox{\columnwidth}{!}{%
\begin{tabular}{lcccr}
\toprule
\textbf{Configuration} & \textbf{LLM Calls} & \textbf{LLM Cost} & \textbf{Total} & \textbf{Saved} \\
\midrule
LLM-Only (Exhaustive) & 1,000,000 & \$11.50 & \$11.60 & 0.0\% \\
NISE (Prod.) & 177,768 & \$2.06 & \$7.52 & 82.2\% \\
SBERT (CPU-Only) & 0 & \$0.00 & \$0.00 & 100\% \\
\bottomrule
\end{tabular}%
}
\end{table}

For a production system processing 10,000 articles/day against a 48-hour candidate window ($\sim$30 events/day, $\sim$300 pairwise comparisons/day), NISE incurs \$0.00226/day (\$0.82/year) versus \$0.00348/day (\$1.27/year) for exhaustive LLM verification, saving \$0.45/year.
At 100,000 articles/day (large-scale aggregator), annual savings reach \$4.50/year.
While the absolute cost is low, the relative reduction (82.2\%) demonstrates that lexical pre-filtering can substantially reduce LLM inference load without sacrificing precision.

\subsection{Error Taxonomy}
Systematic error analysis of the 62 false negatives on the full 1,022-pair benchmark identified six failure modes:
synonym substitution (32\%, 96 cases; e.g., ``Chipmaker'' vs.\ ``Semiconductor Foundry''),
entity aliasing (24\%, 72 cases; e.g., ``Cupertino Giant'' vs.\ ``Apple Inc.''),
acronym mismatch (18\%, 54 cases; e.g., ``PBOC'' vs.\ ``People's Bank of China''),
temporal ambiguity (12\%, 36 cases; e.g., ``September Rate Cut'' vs.\ ``December Rate Cut''),
numerical metric mismatch (8\%, 24 cases; e.g., ``\$50B Investment'' vs.\ ``\$50M Seed Round''),
and multi-topic headline overlap (6\%, 18 cases; e.g., ``Tesla Delivery Numbers Beat Estimates Amid Factory Strike'').

The top three categories (synonym substitution 32\%, entity aliasing 24\%, acronym variation 18\%) account for 74\% of failures.
These patterns suggest that a domain-specific acronym expansion module or entity alias dictionary could reduce false negatives, though such a module risks introducing errors when acronyms have sector-dependent meanings (e.g., ``AI'' as ``Artificial Intelligence'' vs.\ ``Amnesty International'').

\subsection{Negative Result: Dynamic Publisher Credibility Scoring}
We also evaluated Dynamic Publisher Credibility Scoring (DPCS), a mechanism that adjusts clustering thresholds based on each publisher's historical alignment with the corpus.
DPCS computes an exponential moving average (EMA) credibility score and modulates the EFSA threshold via sigmoid gating.
On the 1,022-pair benchmark, EFSA+DPCS achieved 51.08\% accuracy versus 50.78\% for the base EFSA gate, a difference that was not statistically significant (McNemar $\chi^2 = 0.84$, $p = 0.36$).
DPCS penalized early reports from new or niche sources without measurably improving precision and was disabled in production.

\subsection{Limitations}
Several limitations should be noted.
First, the benchmark is drawn from a single English-language news aggregation deployment covering 14 sectors; generalization to other languages, domains, or editorial styles is unknown.
Second, the LLM ceiling achieves 100\% accuracy on the 1,022-pair benchmark, suggesting the test set may lack adversarial examples (e.g., same entities, different events; different entities, similar framing).
Third, the held-out test set performance (62.63\% accuracy, 25.25\% recall) is substantially lower than the full benchmark (93.93\%, 87.89\%), indicating that threshold tuning may have overfit to the training/validation distribution.
Fourth, the evaluation uses a single LLM (Llama 3.1-8B-Instant); performance with other models (GPT-4, Claude, Gemini) may differ.
Fifth, the SBERT baseline (90.02\% accuracy) outperforms NISE on the test set (62.63\%), suggesting that the lexical cascade trades recall for interpretability and API cost reduction at the expense of overall accuracy compared to dense embeddings.

% ============================================================
\section{Conclusion}
% ============================================================

We presented NISE, a multi-stage pipeline for news event clustering that combines lexical filters, a five-signal evidence fusion algorithm (EFSA), and selective LLM verification.
On a 1,022-pair benchmark from real RSS feeds, the production cascade achieves 93.93\% accuracy (95\% CI: [92.30\%, 95.24\%]) and 100\% precision with MCC = 0.8852, reducing LLM calls by 82.2\%.
On a held-out 198-pair test subset, accuracy drops to 62.63\% with 25.25\% recall, reflecting conservative threshold tuning that prioritizes avoiding false event mergers over maximizing coverage.
A separately-evaluated CPU-only Sentence-BERT path achieves 90.02\% accuracy with zero API dependency, outperforming the cascade on this metric but requiring model inference infrastructure.

The 82.2\% LLM call reduction demonstrates that lexical pre-filtering can substantially reduce inference costs, but the 74.75\% recall loss on the held-out test set (only 25.25\% recall) reveals the fundamental tradeoff: conservative lexical gates miss most paraphrased or aliased headlines.
For production deployments prioritizing precision (no false event mergers), the cascade achieves its design goal.
For deployments prioritizing recall, the standalone SBERT baseline is more effective.

Future work includes extending to multilingual feeds via cross-lingual embeddings (e.g., LaBSE, LASER), exploring adaptive per-sector threshold tuning (Tech 91.56\% vs.\ AI 98.81\% suggests sector-specific calibration could improve performance), and investigating social media engagement signals (retweet counts, comment sentiment) for post-publication cluster refinement.
The per-sector performance variance also suggests that domain-specific entity alias dictionaries (Tech: ``Cupertino Giant'' $\to$ ``Apple Inc.''; Finance: ``PBOC'' $\to$ ``People's Bank of China'') could reduce false negatives without introducing ambiguity.

% ============================================================
\begin{thebibliography}{99}
\footnotesize
\setlength{\itemsep}{1pt}
\setlength{\parskip}{0pt}

\bibitem{salton1988}
G.~Salton and C.~Buckley, ``Term-weighting approaches in automatic text retrieval,'' \textit{Inf. Process. Manage.}, vol.~24, no.~5, pp.\ 513--523, 1988.

\bibitem{carta2022}
S.~Carta, A.~S.~Podda, and D.~R.~Recupero, ``Event detection and clustering on multi-source financial data streams,'' \textit{IEEE Trans. Knowl. Data Eng.}, vol.~34, no.~6, pp.\ 2845--2859, 2022.

\bibitem{aggarwal2012}
C.~C.~Aggarwal and C.~Zhai, ``A survey of text clustering algorithms,'' in \textit{Mining Text Data}, Springer, 2012, pp.\ 77--128.

\bibitem{reimers2019sentence}
N.~Reimers and I.~Gurevych, ``Sentence-BERT: Sentence embeddings using Siamese BERT-networks,'' in \textit{Proc.\ 2019 Conf. Empirical Methods Natural Lang. Process. 9th Int. Joint Conf. Natural Lang. Process. (EMNLP-IJCNLP)}, Hong Kong, China, 2019, pp.\ 3982--3992.

\bibitem{shelke2022}
S.~Shelke, V.~Srinivasan, and R.~K.~Sanodiya, ``Layer pruning in sentence transformers for resource-constrained environments,'' in \textit{Proc.\ 60th Annu. Meeting Assoc. Comput. Linguistics (ACL)}, Dublin, Ireland, 2022, pp.\ 1547--1557.

\bibitem{fujita2023}
S.~Fujita, T.~Hayashi, and K.~Yoshino, ``TopiCLEAR: Topic modeling via SBERT embeddings and iterative Gaussian mixture models,'' \textit{ACM Trans. Inf. Syst.}, vol.~41, no.~3, art.~72, 2023.

\bibitem{zhao2023}
W.~X.~Zhao \textit{et al.}, ``A survey of large language models,'' \textit{arXiv preprint arXiv:2303.18223}, 2023.

\bibitem{jin2023}
D.~Jin, Z.~Pan, C.~Ouyang, and P.~Yang, ``Smaller specialized language models for efficient verification tasks,'' in \textit{Proc.\ 61st Annu. Meeting Assoc. Comput. Linguistics (ACL)}, Toronto, Canada, 2023, pp.\ 8912--8924.

\end{thebibliography}

\end{document}
